%%%%%%%%%%%%%%%%%%%%%%%%%%%%%%%%%%%%%%%%%%%%%%%%%%%%%%%%%%%%%%%%%%%%%
%% sec_theory.tex
%% Section 2: Methodology / Theory
%% RMG Periodic RT-TDDFT — Paper 2
%%
%% Subsection structure:
%%   2.1  Overview of RMG (brief, refer to paper 1)
%%   2.2  Interaction with electromagnetic radiation
%%   2.3  Gauge transformation: length to velocity
%%   2.4  Vector potential and magnetic perturbations
%%   2.5  Current operator and nonlocal pseudopotentials  [2d+2e merged]
%%   2.6  Density matrix propagation
%%   2.7  Spin-polarized extension
%%   2.8  Postprocessing: optical conductivity and Berry phase
%%%%%%%%%%%%%%%%%%%%%%%%%%%%%%%%%%%%%%%%%%%%%%%%%%%%%%%%%%%%%%%%%%%%%

\section{Methodology}
\label{sec:theory}

% Units: atomic units (hbar=1, m_e=1, e=1, 4*pi*eps_0=1) throughout
% this section unless stated otherwise.
% Electron charge: -e where e > 0.
% Imaginary unit: \imath (dotless i), consistent with paper 1.

All equations in this section are given in atomic units
($\hbar = m_e = e = 4\pi\varepsilon_0 = 1$) unless stated otherwise.


%--------------------------------------------------------------------
\subsection{Overview of RMG}
\label{sec:rmg-overview}
%--------------------------------------------------------------------

\myred{[Brief overview of RMG: real-space grid, multigrid acceleration,
GPU scalability. Refer heavily to paper 1~\cite{jakowski-rmg-tddft-2025}
and the RMG code paper~\cite{rmg-dft-2024} rather than repeating.
Keep to 1--2 paragraphs. Focus on what is new or relevant for the
periodic case: k-point sampling, Bloch states, periodic boundary
conditions. Mention norm-conserving pseudopotentials (NCPP only).]}

\myred{[Note from CONTEXT.md: the code uses norm-conserving PP only
(asserted in RmgTddft.cpp line 287). Ultrasoft PP requires S$\neq$1
in the propagation (eldyn\_nonort.cpp path) and is not yet implemented.
State this as an explicit scope limitation here.]}


%--------------------------------------------------------------------
\subsection{Interaction of Matter with Electromagnetic Radiation}
\label{sec:em-interaction}
%--------------------------------------------------------------------

\myred{[DERIVATION NEEDED from Jacek's LaTeX files.
Content outline:
- Start from the full minimal-coupling Hamiltonian with EM field
  treated as a plane wave: $\exp(\imath\omega t - \imath\mathbf{q}\cdot\mathbf{r})$
- Derive the long-wavelength (dipole) approximation: $\mathbf{q}\to 0$,
  i.e., $\lambda \gg$ system size, valid for optical/UV and unit-cell scales
- Show that in this limit the vector potential $\mathbf{A}$ becomes
  spatially uniform: $\mathbf{A}(\mathbf{r},t) \to \mathbf{A}(t)$
- Result: the two natural gauge choices in the long-wavelength limit
  are the length gauge (scalar potential $\Phi = -\mathbf{r}\cdot\mathbf{E}$,
  $\mathbf{A}=0$) and the velocity gauge ($\Phi=0$, $\mathbf{A}(t)$ uniform)
- Physical interpretation: length gauge couples to position (dipole),
  velocity gauge couples to momentum (current)
- Reference: Mattiat \& Luber (2022) Section 2.3;
  Yabana \& Bertsch (1996)]}

%% Placeholder equations (to be replaced by Jacek's derivation):

The time-dependent Kohn-Sham (TDKS) equations governing the electronic
dynamics under an external electromagnetic perturbation take the form
\begin{equation}
  \imath \frac{\partial}{\partial t} \psi_n(\mathbf{r},t)
  = \hat{H}_{\mathrm{KS}}(t)\, \psi_n(\mathbf{r},t),
  \label{eq:tdks}
\end{equation}
where $\hat{H}_{\mathrm{KS}}(t)$ is the time-dependent Kohn-Sham
Hamiltonian.
\myred{[Expand: write out $\hat{H}_{\mathrm{KS}}(t)$ with all terms
including the EM coupling, then derive the long-wavelength limit.]}


%--------------------------------------------------------------------
\subsection{Gauge Transformation: Length to Velocity Gauge}
\label{sec:gauge}
%--------------------------------------------------------------------

\myred{[DERIVATION NEEDED from Jacek's LaTeX files.
Content outline:
- State clearly why length gauge fails for periodic systems:
  the position operator $\hat{\mathbf{r}}$ is not translationally
  invariant; dipole moment of a periodic system is multivalued
  (refer to modern theory of polarization: Resta 1994, King-Smith
  \& Vanderbilt 1993)
- Present the G\"{o}ppert-Mayer gauge transformation explicitly:
  gauge function $\chi(\mathbf{r},t) = \mathbf{r}\cdot\mathbf{A}(t)$,
  transformations of potentials and wavefunctions (Mattiat eqs. 15--19)
- Show that the length gauge TDKS (eq.~\ref{eq:tdks-length}) transforms
  to the velocity gauge TDKS (eq.~\ref{eq:tdks-velocity})
- CRITICAL: the nonlocal pseudopotential V\^{}nl does NOT commute with
  the gauge transformation. The gauge-transformed V\^{}nl acquires an
  additional term (Mattiat eq. 23). Derive this explicitly here or
  refer to Appendix~\ref{app:gauge_nlpp}.
- Result: velocity gauge TDKS Hamiltonian with A(t) coupling
- Reference: Mattiat \& Luber (2022) Sections 2.3.1, Appendix A, B;
  Pickard \& Mauri (2001)]}

%% Length gauge TDKS (placeholder):
In the length gauge, the TDKS equations read
\begin{equation}
  \imath \frac{\partial}{\partial t} \psi_n(\mathbf{r},t)
  = \left[ \hat{H}_0 + \mathbf{r}\cdot\mathbf{E}(t) \right]
    \psi_n(\mathbf{r},t),
  \label{eq:tdks-length}
\end{equation}
where $\hat{H}_0$ is the unperturbed Kohn-Sham Hamiltonian and
$\mathbf{E}(t)$ is the electric field.
Applying the gauge transformation $\chi(\mathbf{r},t) =
\mathbf{r}\cdot\mathbf{A}(t)$, where the vector potential satisfies
$\mathbf{E}(t) = -\partial\mathbf{A}/\partial t$,
the velocity gauge TDKS equations take the form
\begin{equation}
  \imath \frac{\partial}{\partial t} \psi_n(\mathbf{r},t)
  = \left[ \frac{1}{2}\left(\hat{\mathbf{p}} + \mathbf{A}(t)\right)^2
    + V(\mathbf{r},t) + \tilde{V}^{\mathrm{nl}}(\mathbf{r},t)
  \right] \psi_n(\mathbf{r},t),
  \label{eq:tdks-velocity}
\end{equation}
where $\tilde{V}^{\mathrm{nl}}$ denotes the gauge-transformed nonlocal
pseudopotential.\cite{Mattiat-Luber-2022,pickard-mauri-2001}
\myred{[Expand: write out all terms of $\tilde{V}^{\mathrm{nl}}$
explicitly. Derive the linearization in A(t) that gives the
perturbative coupling term. Connect to Appendix~\ref{app:gauge_nlpp}.]}


%--------------------------------------------------------------------
\subsection{Vector Potential and Magnetic Perturbations}
\label{sec:vector-potential}
%--------------------------------------------------------------------

\myred{[DERIVATION NEEDED from Jacek's LaTeX files.
Content outline:
- Show that the velocity gauge framework with $\mathbf{A}(t)$
  naturally accommodates both electric and magnetic perturbations
- Electric perturbation in the long-wavelength limit:
  $\mathbf{A}(t)$ spatially uniform, $\mathbf{B}=0$
- Magnetic perturbation (symmetric gauge):
  $\mathbf{A}(\mathbf{r}) = \frac{1}{2}\mathbf{B}\times(\mathbf{r}-\mathbf{R}_G)$,
  spatially dependent, introduces gauge origin $\mathbf{R}_G$
- Note: magnetic perturbation sets up ECD/MCD as future extension
  (not implemented in present work but theoretically grounded here)
- The gauge origin dependence with finite basis sets and the standard
  remedies (GIAO, IGLO) -- brief mention, refer to Mattiat for details
- Reference: Mattiat \& Luber (2022) Section 2.3.2;
  Pickard \& Mauri (2001)]}

\myred{[Note: This subsection establishes the theoretical foundation
for magnetic perturbations. The current paper implements only electric
perturbations (A(t) uniform). ECD/MCD calculations are future work.
State this explicitly at the end of the subsection.]}


%--------------------------------------------------------------------
\subsection{Current Operator and Nonlocal Pseudopotential Contributions}
\label{sec:current}
%--------------------------------------------------------------------

\myred{[DERIVATION NEEDED from Jacek's LaTeX files.
This is the most important new subsection. Full derivation required.
Content outline:

(a) Canonical vs mechanical momentum:
    - Canonical: $\hat{\mathbf{p}} = -\imath\nabla$
    - Mechanical: $\hat{\boldsymbol{\pi}} = \hat{\mathbf{p}} + \mathbf{A}(t)$
    - Velocity operator: $\hat{\mathbf{v}} = (1/\imath)[\hat{\mathbf{r}}, \hat{H}]$
    - When $V^{\mathrm{nl}} \neq 0$: $\hat{\mathbf{v}} \neq \hat{\mathbf{p}}/m$
      -- extra commutator term $[\hat{\mathbf{r}}, V^{\mathrm{nl}}]$
    - Ref: Starace (1971,1973); Mattiat \& Luber eq.12

(b) Paramagnetic and diamagnetic current:
    - Full current: $\hat{\mathbf{J}} = \hat{\mathbf{J}}_{\mathrm{para}} + \hat{\mathbf{J}}_{\mathrm{dia}}$
    - Paramagnetic: $\hat{\mathbf{J}}_{\mathrm{para}} \propto \mathrm{Re}[\psi^*(-\imath\nabla)\psi]$
    - Diamagnetic: $\hat{\mathbf{J}}_{\mathrm{dia}} \propto \mathbf{A}(t)|\psi|^2$
    - J\_dia = 0 after t=0 in our perturbative scheme (A=0 for t>0)
    - Justify this explicitly as valid in linear response regime

(c) Continuity equation:
    - $\partial\rho/\partial t + \nabla\cdot\mathbf{J} = 0$
    - Use as consistency/sanity check for the current operator

(d) Current in density matrix formulation:
    - Macroscopic current: $J_\alpha(t) = \mathrm{Re}\,\mathrm{Tr}[\mathbf{P}(t)\cdot\mathbf{J}^\alpha]$
    - Current operator matrix: $J^\alpha_{ab} = p^\alpha_{ab} + [\hat{r}_\alpha, V^{\mathrm{nl}}]_{ab}$
    - where $p^\alpha_{ab} = \langle\psi_a|(-\imath\nabla_\alpha + k_\alpha)|\psi_b\rangle$
      (the $k_\alpha$ term comes from the Bloch theorem: see eq.XX)
    - This is what VecPHmatrix() + CurrentNlpp() compute in the code

(e) Nonlocal PP contribution [r, V^nl]:
    - For Kleinman-Bylander separable pseudopotentials:
      $V^{\mathrm{nl}} = \sum_{I,lm} |\beta^{lm}_I\rangle h^l_{ij} \langle\beta^{lm}_I|$
    - Compute $[\hat{r}_\alpha, V^{\mathrm{nl}}]$ explicitly:
      $[\hat{r}_\alpha, V^{\mathrm{nl}}]|\psi\rangle =
       \sum_{I,lm} |\beta^{lm}_I\rangle h^l_{ij}
       \langle[\hat{r}_\alpha, \beta^{lm}_I]|\psi\rangle$
    - The commutator acts on the projector: $[\hat{r}_\alpha, \beta^{lm}_I]$
      is computed by AppNls\_0xyz() in the code
    - Extended algebra in Appendix~\ref{app:gauge_nlpp}
    - Ref: Mattiat \& Luber (2022) Section 2.2, eq.12;
           Pickard \& Mauri (2001)

VERIFICATION NEEDED: trace the sign/i-factor from theory to code.
VecPHmatrix() stores I\_t * block\_matrix (I\_t = i) in Pxyz.
CurrentNlpp() adds I\_t * block\_matrix.
Current extracted as Re(zdotc(P, Px)).
The -i from the gradient operator + i from storage = real observable.
Verify this chain explicitly in the derivation.
See Gap 1 in CONTEXT.md.]}

%% Placeholder content (to be replaced by full derivation):

The velocity operator for a system described by a Hamiltonian
containing a nonlocal pseudopotential $\Vnl$ is obtained from the
Heisenberg equation of motion,
\begin{equation}
  \hat{v}_\alpha = \frac{1}{\imath}[\hat{r}_\alpha, \hat{H}]
  = \hat{p}_\alpha + \frac{1}{\imath}[\hat{r}_\alpha, \Vnl],
  \label{eq:velocity-op}
\end{equation}
where the second term on the right-hand side arises because
$[\hat{r}_\alpha, \Vnl] \neq 0$ for a nonlocal
operator.\cite{starace-1971,Mattiat-Luber-2022}
\myred{[Expand into full derivation of current operator, paramagnetic
and diamagnetic decomposition, density matrix formulation, and explicit
KB pseudopotential commutator. Follow derivation philosophy: complete
first, then decide what moves to appendix.]}

The macroscopic current density is expressed in terms of the
one-particle density matrix $\vb{P}(t)$ as
\begin{equation}
  J_\alpha(t) = \mathrm{Re}\,\mathrm{Tr}\!\left[
    \vb{P}(t)\cdot\vb{J}^\alpha \right],
  \label{eq:current-dm}
\end{equation}
where $\vb{J}^\alpha$ is the matrix of current operator elements
in the Kohn-Sham orbital basis,
\begin{equation}
  J^\alpha_{ab} = \langle\psi_a|
    \left(-\imath\nabla_\alpha + k_\alpha
    + \frac{1}{\imath}[\hat{r}_\alpha, \Vnl]\right)
  |\psi_b\rangle.
  \label{eq:current-matrix}
\end{equation}
\myred{[The $k_\alpha$ term in eq.~\ref{eq:current-matrix} arises from
the Bloch theorem: for a periodic orbital
$\psi_{n,\mathbf{k}}(\mathbf{r}) = e^{\imath\mathbf{k}\cdot\mathbf{r}}
u_{n,\mathbf{k}}(\mathbf{r})$,
the gradient operator acting on the full Bloch function gives
$\nabla(e^{\imath\mathbf{k}\cdot\mathbf{r}}u) =
e^{\imath\mathbf{k}\cdot\mathbf{r}}(\nabla + \imath\mathbf{k})u$.
Derive this step explicitly.]}


%--------------------------------------------------------------------
\subsection{Density Matrix Propagation}
\label{sec:propagation}
%--------------------------------------------------------------------

\myred{[Adapt from paper 1 (Jakowski et al. 2025, Section 2.2).
The derivation is identical to paper 1 for the case S=1 (orthogonal
basis at Gamma point). For k \neq Gamma, H is complex-valued (Hermitian
but not real-symmetric) and Omega is also complex. The BCH expansion
is the same but with complex matrix arithmetic.
Key additions vs paper 1:
- Note explicitly that for k \neq Gamma the Hamiltonian matrix
  $\mathbf{H}_\mathbf{k} = \langle\psi_a|\hat{H}|\psi_b\rangle$
  is complex Hermitian
- The BCH propagator (Ieldyn=1) handles complex matrices;
  the diagonalization path (Ieldyn=2) does NOT -- state this
- The self-consistency loop (predictor-corrector) is unchanged
- The Hamiltonian at each time step includes the A(t)*p coupling
  term (for VECTOR\_POT mode); in the perturbative kick scheme
  A appears only in the initial H\_0 setup
Reference sections from paper 1: Magnus expansion and density matrix
propagation subsections.
Cite: Jakowski \& Morokuma (2009); Jakowski et al. JCTC 2025.]}

The time evolution of the one-particle density matrix $\vb{P}(t)$
is governed by the Liouville-von Neumann equation
\begin{equation}
  \frac{\partial \vb{P}(t)}{\partial t}
  = -\imath\left[\vb{H}(t),\, \vb{P}(t)\right],
  \label{eq:LvN}
\end{equation}
where $\vb{H}(t)$ is the Kohn-Sham Hamiltonian matrix in the basis
of ground-state Kohn-Sham orbitals.
\myred{[Continue with Magnus expansion (refer to paper 1 eqs. and
cite), then commutator expansion for P(t+dt), self-consistency loop.
For complex H: add note about BCH-only path.]}

\myred{[Note from code inspection (CONTEXT.md Gap 4): magnus() in the
code is documented as real-only, but for k$\neq\Gamma$ the call site
casts to double*. This is mathematically correct since the Magnus
operation $\Omega = -\imath\cdot\frac{1}{2}(H_0+H_1)\Delta t$ is
element-wise and valid for both real and complex H, but should be
stated explicitly in the text.]}


%--------------------------------------------------------------------
\subsection{Spin-Polarized Extension}
\label{sec:spin}
%--------------------------------------------------------------------

\myred{[DERIVATION NEEDED from Jacek's LaTeX files.
Content outline:
- Collinear spin-polarized DFT: two independent sets of KS orbitals
  $\{\psi^\uparrow_n\}$ and $\{\psi^\downarrow_n\}$, each with its
  own density matrix $\vb{P}^\uparrow(t)$ and $\vb{P}^\downarrow(t)$
- Spin-resolved electron density:
  $\rho^\uparrow(\mathbf{r},t)$ and $\rho^\downarrow(\mathbf{r},t)$
- Total density: $\rho = \rho^\uparrow + \rho^\downarrow$
- Spin magnetization: $m_s = \rho^\uparrow - \rho^\downarrow$
- Coupling between spin channels: through the spin-dependent XC potential
  $V_{\mathrm{xc}}^\sigma[\rho^\uparrow, \rho^\downarrow]$
  (LSDA or GGA spin-polarized functional)
- Each spin channel propagates independently:
  $\partial\vb{P}^\sigma/\partial t = -\imath[\vb{H}^\sigma(t), \vb{P}^\sigma(t)]$
- The Hamiltonians $\vb{H}^\uparrow$ and $\vb{H}^\downarrow$ differ
  only through $V_{\mathrm{xc}}^\uparrow \neq V_{\mathrm{xc}}^\downarrow$
- Spin-resolved current and optical conductivity:
  $J^\sigma_\alpha(t) = \mathrm{Re}\,\mathrm{Tr}[\vb{P}^\sigma(t)\cdot\vb{J}^\alpha]$
  $\boldsymbol{\sigma}^\sigma(\omega)$ for each spin channel
- Total and spin currents:
  $\mathbf{J}_{\mathrm{charge}} = \mathbf{J}^\uparrow + \mathbf{J}^\downarrow$
  $\mathbf{J}_{\mathrm{spin}} = \mathbf{J}^\uparrow - \mathbf{J}^\downarrow$
- Note: non-collinear spin (two-component spinors) is beyond the scope
  of this work and will be presented in a future publication.
- Note from code inspection: spin is handled via separate MPI
  communicator groups (pct.spinpe); spin channels communicate only
  when updating $V_{\mathrm{xc}}$ (rho.get\_oppo()).]}

In the collinear spin-polarized extension, the density matrix
propagation (Section~\ref{sec:propagation}) is performed independently
for each spin channel $\sigma \in \{\uparrow, \downarrow\}$,
\begin{equation}
  \frac{\partial \vb{P}^\sigma(t)}{\partial t}
  = -\imath\left[\vb{H}^\sigma(t),\, \vb{P}^\sigma(t)\right],
  \label{eq:LvN-spin}
\end{equation}
where the spin-dependent Hamiltonian $\vb{H}^\sigma(t)$ contains the
spin-resolved exchange-correlation potential
$V^\sigma_{\mathrm{xc}}[\rho^\uparrow(\mathbf{r},t),
\rho^\downarrow(\mathbf{r},t)]$.
\myred{[Continue: derive spin-resolved current, define total and spin
currents, note that J\_spin = J\_up - J\_down is available as a
low-cost observable.]}


%--------------------------------------------------------------------
\subsection{Postprocessing: Optical Conductivity and Berry Phase}
\label{sec:postproc}
%--------------------------------------------------------------------

\myred{[Content outline:
(a) From J(t) to optical conductivity:
    - Fourier transform: $\tilde{J}_\alpha(\omega) = \int J_\alpha(t) e^{\imath\omega t} dt$
    - Optical conductivity tensor:
      $\sigma_{\alpha\beta}(\omega) = \tilde{J}_\alpha(\omega) / E_\beta(\omega)$
    - Dielectric function:
      $\varepsilon_{\alpha\beta}(\omega) = \delta_{\alpha\beta} +
       \imath\sigma_{\alpha\beta}(\omega)/\varepsilon_0\omega$
    - In practice: apply delta-kick perturbation in direction $\beta$,
      Fourier transform $J_\alpha(t)$, divide by $E_\beta(\omega)=E_0$
      (constant for delta kick)
    - Mention damping/broadening: apply $e^{-\gamma t}$ window before FFT

(b) Berry phase polarization as alternative observable:
    - For periodic systems, polarization can also be extracted via the
      King-Smith/Vanderbilt Berry phase formula:\cite{king-smith-vanderbilt-1993}
      $\mathbf{P}_{\mathrm{BP}} = \frac{e}{(2\pi)^3}\int_{\mathrm{BZ}}
       \mathbf{A}_n(\mathbf{k})\,d\mathbf{k}$
      where $\mathbf{A}_n = \imath\langle u_{n\mathbf{k}}|\nabla_\mathbf{k}|u_{n\mathbf{k}}\rangle$
      is the Berry connection
    - For the Gamma-point supercell limit, the discrete formula
      (Resta 1998, eq. 48 in Mattiat \& Luber):
      $\langle\hat{d}_\alpha\rangle = \frac{eL}{2\pi}\mathrm{Im}\ln\det S_\alpha$
      where $S_{\alpha,ij} = \langle\psi_i|e^{-2\pi\imath r_\alpha/L}|\psi_j\rangle$
    - In density matrix form: $P_{\mathrm{BP},\alpha}(t) = \mathrm{Tr}[\vb{P}(t)\cdot\vb{X}^\alpha]$
      where $\vb{X}^\alpha$ is the Berry phase position matrix
    - Equivalence with current-integration result:
      $\Delta P_\alpha = \int_0^t J_\alpha(t')\,dt'$
      (test of internal consistency)
    - Implementation in RMG: BP\_Xml matrix computed by
      Rmg\_BP->tddft\_Xml() (see code, RmgTddft.cpp lines 514,517)
    - Reference: King-Smith \& Vanderbilt (1993); Resta (1998);
      Mattiat \& Luber (2022) Section 2.5]}

The macroscopic current density $\mathbf{J}(t)$ obtained from
eq.~(\ref{eq:current-dm}) provides direct access to the optical
response of the periodic system.
After applying a momentum kick of amplitude $E_0$ along the Cartesian
direction $\beta$ at $t=0$, the optical conductivity tensor is obtained
from the Fourier transform of the current density as
\begin{equation}
  \sigma_{\alpha\beta}(\omega)
  = \frac{\tilde{J}_\alpha(\omega)}{E_\beta(\omega)}
  = \frac{1}{E_0}\int_0^\infty J_\alpha(t)\,e^{\imath\omega t - \gamma t}\,dt,
  \label{eq:conductivity}
\end{equation}
where $\gamma$ is a broadening parameter.\cite{yabana-bertsch-1996}
\myred{[Expand: relate $\sigma$ to $\varepsilon$; describe how three
independent kick simulations (x, y, z) give the full tensor; note
crystal symmetry reduces the number of independent components.]}

\myred{[Add: Berry phase derivation and density matrix form.
Show equivalence with current integration.
Note that both observables are computed in the present implementation
and their agreement serves as an internal consistency check.]}
